\chapter{Hawking Radiation}
\label{ch:hawking-radiation}
We now bring together all the tools developed in this part to derive
the central result: a black hole formed by gravitational collapse emits
a steady thermal flux of radiation at the Hawking temperature
$T_H = 1/8\pi M$. The derivation proceeds by tracing the relation
between modes defined on past null infinity and modes defined on future
null infinity, through the collapsing geometry. The logarithmic
relationship between the in and out coordinates near the horizon, which
we already encountered in the accelerating mirror, is the key
ingredient. The result connects back to the thermodynamic picture of
Chapter~\ref{ch:bh-thermodynamics}: the Hawking temperature is
precisely the temperature one would assign to a black hole of entropy
$S_{BH} = A/4$ using the thermodynamic relation $T = dM/dS$. It also
connects forward to the cosmological context of Part~\ref{part:cosmo}:
the same mechanism, quantum fluctuations near a causal horizon, is
responsible for the Gibbons-Hawking temperature of de Sitter space and
ultimately for the primordial density perturbations seeded during
inflation.

\section{Particle Production in Gravitational Collapse}

We now derive the most important consequence of quantum field theory
in curved spacetime: the thermal radiation emitted by a black hole,
discovered by Hawking in 1974. We model the gravitational collapse by
a thin spherical shell of mass $M$ collapsing to form a black hole.
By Birkhoff's theorem, the metric outside the shell is the
Schwarzschild metric,
\begin{equation}
    ds^2_{\text{out}} = -f(r)\,du\,dv,
    \qquad f(r) = 1 - \frac{2M}{r},
    \qquad u,v = t\mp r_*,
\end{equation}
and the metric inside the shell is flat,
\begin{equation}
    ds^2_{\text{in}} = -d\tilde u\,d\tilde v,
    \qquad \tilde u,\tilde v = \tilde t\mp r.
\end{equation}

% [FIGURE: Penrose diagram of gravitational collapse with thin shell,
%  showing past null infinity, future null infinity, the horizon,
%  and the singularity]

The history of the collapsing matter does not affect the particle
production; what matters is the relationship between the coordinates
on either side of the shell near the moment of horizon crossing.

\subsection*{The 2D Model}

We work in the 2D reduction, suppressing the angular coordinates. The
shell trajectory is $r = r_s(\tilde t)$, and matching the metrics
across the shell gives
\begin{align}
    \frac{du}{d\tilde u}
    &= \frac{\sqrt{(1-\dot r_s^2)f(r_s)+\dot r_s^2}-\dot r_s}
    {f(r_s)(1-\dot r_s)}, \\
    \frac{dv}{d\tilde v}
    &= \frac{\sqrt{(1-\dot r_s^2)f(r_s)+\dot r_s^2}+\dot r_s}
    {f(r_s)(1+\dot r_s)},
\end{align}
where $\dot r_s = dr_s/d\tilde t$. As the shell crosses the horizon,
$f(r_s)\to 0$. Writing $r_s(\tilde t) = r_h + A(\tilde t - \tilde t_h)
+ \cdots$ with $0 > A > -1$ and $f(r) \simeq (r-r_h)/\kappa_h$ near
the horizon, one finds that $dv/d\tilde v$ remains regular at horizon
crossing,
\begin{equation}
    v = \text{const} + \alpha\tilde v\big|_{\text{horizon crossing}},
\end{equation}
so $v$ is a smooth function of $\tilde v$ there. By contrast,
$du/d\tilde u$ diverges logarithmically,
\begin{equation}
    u = -\kappa_h\ln|\tilde u - \tilde u_h| + \text{const},
\end{equation}
where $\kappa_h = f'(r_h)/2$ is the surface gravity. This logarithmic
relation is the key to the thermal spectrum.

The natural in-modes on past null infinity $\mathscr{I}^-$ are
\begin{equation}
    \varphi^{\text{in}}_\omega
    = \frac{1}{\sqrt{4\pi\omega}}
    \left(e^{-i\omega v} - e^{-i\omega u}\right),
\end{equation}
and the natural out-modes for an asymptotic observer on future null
infinity $\mathscr{I}^+$ are
\begin{equation}
    \varphi^{\text{out}}_{\tilde\omega}
    = \frac{1}{\sqrt{4\pi\tilde\omega}}
    \left(e^{-i\tilde\omega u}
    - e^{-i\tilde\omega u_0(v)}\right),
\end{equation}
where $u_0(v)$ is the image of the out-ray under reflection from the
center. Using the logarithmic relation,
\begin{equation}
    u_0 = -\kappa_h\ln|v - v_h| + \text{const},
\end{equation}
so the out-modes contain the same logarithmic phase as the accelerating
mirror. The computation of the Bogoliubov coefficients therefore
proceeds identically to the mirror case with the identification
$K\to\kappa_h$, giving a thermal flux at temperature
\begin{equation}
    \boxed{T_H = \frac{\kappa_h}{2\pi}
    = \frac{f'(r_h)}{4\pi}
    = \frac{1}{8\pi M}.}
\end{equation}
This is the Hawking temperature. The particle production rate per unit
retarded time is
\begin{equation}
    \frac{dN}{d\tilde\omega\,du}
    = \frac{1}{2\pi}\cdot\frac{1}{e^{\tilde\omega/T_H}-1},
\end{equation}
a steady thermal flux that persists as long as the black hole exists.
An observer at $\mathscr{I}^+$ will keep receiving this radiation even
though no signal can escape from inside the horizon; the radiation
originates from quantum fluctuations just outside the horizon, not
from inside it.

\subsection*{The 4D Model}

In four dimensions the analysis is more involved due to the angular
structure. The field modes take the form
\begin{equation}
    \varphi_{\omega lm}
    = \frac{Y_{lm}(\theta,\phi)}{r}\,R_{\omega l}(r)\,e^{-i\omega t},
\end{equation}
where the radial function satisfies a Schr\"{o}dinger-like equation,
\begin{equation}
    -\frac{d^2R_{\omega l}}{dr_*^2}
    + V_{\text{eff}}(r_*)\,R_{\omega l} = \omega^2 R_{\omega l},
    \qquad
    V_{\text{eff}}(r_*) = \left(\frac{l(l+1)}{r^2}
    + \frac{2M}{r^3}\right)\left(1-\frac{2M}{r}\right).
\end{equation}

% [FIGURE: effective potential V_eff(r_*) showing the barrier between
%  the horizon and infinity]

The potential vanishes both at the horizon ($r_*\to-\infty$) and at
infinity ($r_*\to+\infty$), so the modes are free waves in both
asymptotic regions. A wave incident from infinity partially reflects
off the potential barrier and partially transmits through to the
horizon. The transmission coefficient $\Gamma_{\tilde\omega l}$,
called the greybody factor, modifies the thermal flux,
\begin{equation}
    \frac{dN}{d\tilde\omega\,dt}
    = \frac{\Gamma_{\tilde\omega l}}{2\pi}
    \cdot\frac{1}{e^{\tilde\omega/T_H}-1}.
\end{equation}
The greybody factor satisfies $\Gamma_{\tilde\omega l}\to 0$ as
$l\to\infty$ or as $\tilde\omega\to 0$ at fixed $l$, so low-frequency
and high-angular-momentum modes are suppressed. The radiation is
therefore not exactly thermal, but it is in detailed balance with a
thermal bath at temperature $T_H$: the emission and absorption
coefficients are equal by time-reversal symmetry, so a black hole in
a thermal bath at $T_H$ would be in equilibrium.

\section{The Eternal Black Hole and the Choice of Vacuum}

For the eternal (maximally extended) Schwarzschild black hole, the
metric in Kruskal coordinates is
\begin{equation}
    ds^2 = -\frac{2M}{r}e^{-r/2M}\,d\bar u\,d\bar v,
    \qquad
    \bar u = -4Me^{-u/4M},\quad
    \bar v = 4Me^{v/4M}.
\end{equation}
The physical content of the theory depends crucially on the choice of
quantum vacuum state, and there are three natural choices.

\subsection*{The Boulware Vacuum}

The Boulware vacuum $|0\rangle_B$ is defined by the modes
\begin{equation}
    \varphi^B_{\omega,R} = \frac{1}{\sqrt{4\pi\omega}}e^{-i\omega u},
    \qquad
    \varphi^B_{\omega,L} = \frac{1}{\sqrt{4\pi\omega}}e^{-i\omega v},
\end{equation}
which are plane waves for an asymptotically distant observer and
correspond to the natural vacuum seen by a static observer far from
the black hole. However, these modes are defined only in region I of
the Kruskal diagram and are singular at the horizon. The Wightman
function in the Boulware vacuum is
\begin{equation}
    G^+_B = -\frac{1}{4\pi}\ln\left[
    (u-u'-i\varepsilon)(v-v'-i\varepsilon)\right],
\end{equation}
which has no poles in the lower half plane. A static detector at fixed
$r$ therefore measures zero particle flux in the Boulware vacuum,
$F_B(\Delta E) = 0$.

\subsection*{The Hartle-Hawking Vacuum}

The Hartle-Hawking vacuum $|0\rangle_H$ is defined by modes that are
regular everywhere except at the singularity,
\begin{equation}
    \varphi^H_{\omega,R} = \frac{1}{\sqrt{4\pi\omega}}e^{-i\omega\bar u},
    \qquad
    \varphi^H_{\omega,L} = \frac{1}{\sqrt{4\pi\omega}}e^{-i\omega\bar v}.
\end{equation}
These modes are smooth across the horizon and provide a regular
quantization for the full extended spacetime. The Bogoliubov
coefficients between the Hartle-Hawking and Boulware modes are
\begin{equation}
    \beta_{\tilde\omega\omega,RR}
    = -\langle\varphi^B_{\tilde\omega R},
    (\varphi^H_{\omega,R})^*\rangle
    = \frac{-i}{2\pi\sqrt{\omega\tilde\omega}}
    \left(\frac{i}{4M\omega}\right)^{i4M\tilde\omega}
    \Gamma(1+i4M\tilde\omega),
\end{equation}
which are nonzero and describe a thermal flux between the black hole
and its environment. The Wightman function is
\begin{equation}
    G^+_H(t,r,t',r')
    = -\frac{1}{4\pi}\ln\left[
    (\bar u-\bar u'-i\varepsilon)
    (\bar v-\bar v'-i\varepsilon)\right].
\end{equation}
For a static detector at fixed $r$, substituting the trajectory gives
\begin{align}
    G^+_H(t,r,t',r)
    &= -\frac{1}{4\pi}\ln\!\left[
    \cosh\!\left(\frac{t-t'-i\varepsilon}{4M}\right) - 1\right]
    + \tilde G(r) \notag\\
    &= -\frac{1}{4\pi}\ln\!\left[
    \cosh\!\left(\frac{\tau-\tau'-i\varepsilon}
    {4M\sqrt{1-\frac{2M}{r}}}\right) - 1\right]
    + \tilde G(r),
\end{align}
where we used $d\tau = \sqrt{1-2M/r}\,dt$. The Wightman function has
periodic poles in the imaginary time direction with period
\begin{equation}
    \Delta\tau = i\cdot 8\pi M\sqrt{1-\frac{2M}{r}},
\end{equation}
so the detector measures a temperature
\begin{equation}
    T = \frac{1}{8\pi M\sqrt{1-\frac{2M}{r}}}.
\end{equation}
An observer closer to the horizon measures a higher temperature due to
the blueshift factor $\sqrt{1-2M/r}^{-1}$, diverging at the horizon
itself. This is consistent with the Unruh effect: a static observer in
Schwarzschild spacetime is accelerated with respect to a freely falling
frame, and the acceleration grows without bound as the observer
approaches the horizon.

For a freely falling observer crossing the horizon at fixed Kruskal
coordinate $\bar r = (\bar v - \bar u)/2$, the Wightman function is
\begin{equation}
    G^+_H(\bar t,\bar r,\bar t',\bar r)
    = -\frac{1}{2\pi}\ln(\bar t-\bar t'-i\varepsilon),
\end{equation}
which has no poles in the lower half plane, so the freely falling
observer's detector measures nothing. The Hartle-Hawking vacuum is
regular from the perspective of a freely falling observer crossing the
horizon. This state physically describes a black hole in thermal
equilibrium with its own radiation.

\subsection*{The Unruh Vacuum}

For a realistic gravitational collapse, there is no incoming radiation
from past null infinity, only outgoing Hawking radiation. The
appropriate state is the Unruh vacuum $|0\rangle_U$, defined by
\begin{equation}
    \varphi^U_{\omega,R}
    = \frac{1}{\sqrt{4\pi\omega}}e^{-i\omega\bar u},
    \qquad
    \varphi^U_{\omega,L}
    = \frac{1}{\sqrt{4\pi\omega}}e^{-i\omega v}.
\end{equation}
This state uses the Kruskal mode for the right-movers (regular at the
future horizon) and the Boulware mode for the left-movers (no incoming
flux from the past). It is regular at the future horizon but singular
at the past horizon, which is not physically relevant for a collapsing
star. A static detector at fixed $r$ in the Unruh vacuum sees a
one-sided thermal flux of outgoing radiation at the Hawking temperature
$T_H = 1/8\pi M$, with no incoming flux from infinity. This is the
state that correctly describes the endpoint of gravitational collapse.

\section{Hawking Radiation in the Bigger Picture}

The three results we have obtained, the Unruh effect, the thermal
spectrum from the accelerating mirror, and Hawking radiation from a
black hole, are all manifestations of the same underlying physics: the
observer-dependence of the vacuum state in the presence of a causal
horizon. In each case, a state that appears to be the vacuum for one
class of observers contains a thermal distribution of particles for
another class.

For the black hole, the connection is made precise by the
identification of the surface gravity $\kappa_h$ with the acceleration
of a static observer near the horizon. As $r\to r_h$, the proper
acceleration required to remain stationary diverges as
$a\sim 1/(4M\sqrt{1-2M/r})$, and the local Unruh temperature
$T = a/2\pi$ coincides with the blueshifted Hawking temperature
$T_H/\sqrt{1-2M/r}$. The Hawking effect and the Unruh effect are
therefore the same phenomenon seen from different perspectives: global
versus local.

The Hawking temperature $T_H = 1/8\pi M$ implies that black holes
evaporate by radiating energy, causing $M$ to decrease and $T_H$ to
increase. This leads to the information paradox: if the radiation is
exactly thermal, then no information about the initial state that
formed the black hole is encoded in the outgoing radiation, apparently
violating unitarity. The resolution of this paradox remains one of the
central open problems in theoretical physics, at the intersection of
quantum mechanics, general relativity, and thermodynamics.
